% Part: methods
% Chapter: proofs
% Section: starting-proofs

\documentclass[../../../include/open-logic-section]{subfiles}

\begin{document}

\olfileid{mth}{prf}{str}

\olsection{સાબિતીની શરૂઆત}

પરંતુ શરૂઆત ક્યાંથી કરવી?

તમને કંઈક સાબિત કરવા અપાયું છે, તેથી સાબિતીમાં એ બાબત છેલ્લે
આવવી જોઈએ (અલબત્ત, શરૂઆતમાં તમે તેને સાબિત કરવાની \emph{જાહેરાત}
કરી શકો છો, પરંતુ તેને પૂર્વધારણા તરીકે વાપરવી નહીં). તમે જે
સાબિત કરવા માગો છો તે કાગળના નવા પાનાના તળિયે લખો---આ રીતે
તમારું લક્ષ્ય નજર સામે રહેશે.

પછી, કદાચ તમારી પાસે વાપરી શકાય એવી કેટલીક પૂર્વધારણાઓ હશે
(આગળના વિભાગમાં તમારી સાબિતીના \emph{પ્રકાર}ની વાત કરીશું ત્યારે
આ વધુ સ્પષ્ટ થશે). તેને પાનાની ટોચે લખો અને તે પૂર્વધારણાઓ છે
એવું સ્પષ્ટ જણાવો (એટલે, જો તમે $p$ ધારો છો, તો ``ધારો કે $p$''
અથવા ``માની લો કે $p$'' લખો). છેવટે, પ્રશ્નમાં કેટલીક એવી
વ્યાખ્યાઓ હોઈ શકે જે તમારે જાણવી જરૂરી છે. તમને કોઈ ચોક્કસ
વ્યાખ્યા વાપરવાનું કહેવામાં આવ્યું હોઈ શકે, અથવા પૂર્વધારણાઓમાં
કે તમે મેળવવા માગતા નિષ્કર્ષમાં વિવિધ વ્યાખ્યાઓ આવી શકે.
\emph{તેને લખી લો અને તેમનો અર્થ તમને સમજાયો છે તેની ખાતરી કરો.}

તમે સાબિતી કેવી રીતે ગોઠવો તે પ્રશ્નના સ્વરૂપ પર પણ આધાર રાખશે.
આગળનો વિભાગ વાક્યના પ્રકાર પ્રમાણે સાબિતી કેવી રીતે ગોઠવવી તેની
વિગતો આપે છે.

\end{document}
